% MANUSCRIPT STATUS: DIGITAL_RESULTS_VALIDATED
% Evidence class: SYNTHETIC_SIM. Not camera-ready. RESULT_PENDING for any RF/field number.
% Never SUBMITTED / ACCEPTED from this file.
\documentclass[11pt]{article}
\usepackage[margin=1in]{geometry}
\usepackage{booktabs,hyperref,graphicx}
\title{Minimum-Useful-Service Continuity Benchmarking for Resource-Constrained Edge Devices in Heterogeneous Networks}
\author{Edmund Gunn Jr.\\Independent researcher --- no University of Oulu affiliation claimed}
\date{\today}
\begin{document}
\maketitle
\begin{center}\textbf{RESULT BANNER.} Digital synthetic-fixture results only. Independent human reproduction pending. No measured RF. Not SUBMITTED / ACCEPTED.\end{center}
\begin{abstract}
A radio link being available does not imply that a human-facing service remains above its minimum-useful operating point.
This Paper~I manuscript turns four research device classes (desk, mobile/docked, local-creation, wearable) and the Gary flagship scenario into a runnable continuity benchmark.
All numbers are produced by \texttt{make paper-reproduce} from frozen protocol \texttt{paper/artifacts/experiment\_protocol.yaml} and are labelled \texttt{SYNTHETIC\_SIM}.
They are not campus deployments, not RF sounding, and not University of Oulu measurements.
The benchmark's load-bearing result is not that a harness exists: it is \emph{which} class/workload/scenario drops below min-useful, and which labeled stresses do not.
\end{abstract}
\section{Introduction}
A radio link being available does not imply that a human-facing service remains above its minimum-useful operating point~\cite{3gpp_ts22261,itur_m2160}.
Paper~I asks how workloads and four device classes (sustained desk, mobile/docked, local create/deploy, wearable sensing) become measurable continuity profiles and benchmark scenarios.
The primary repositories are \texttt{gunnchos-device-os} and this digital twin.
Gary is the flagship \emph{research scenario}, not a deployed community network.

\section{Method}
Experiment \texttt{rq1\_gary\_flagship\_profiles} is pre-registered in \texttt{configs/experiments/rq1\_gary\_flagship\_profiles.yaml} (seeds, metrics, non-claims) before any held-out RF campaign --- none is claimed here.
Mode is \texttt{synthetic-fixture}. Baselines named for later campaigns: terrestrial-only, cloud-only, local-only, edge-only, fixed policy, no adaptation, oracle. This digital run exports campus metric families (workload, compute, radio, failure, mobility, inclusion) with explicit \texttt{synthetic\_fixture} evidence flags.
Device-class constraints are imported from the device-os continuity profile export when present; otherwise the twin site YAML is used alone.

\section{Results}
\label{sec:results}
Tables in \texttt{paper/tables/} are generated by \texttt{paper/scripts/generate\_tables.py} from \texttt{results/experiments/rq1\_gary\_flagship\_profiles.json}.
If that artifact is missing, the generated TeX file states \texttt{RESULT\_PENDING} rather than inventing numbers.
All present numbers are \texttt{SYNTHETIC\_SIM}.
\begin{table}[h]
\centering
\caption{RQ1 Gary synthetic-user demand by seed (SYNTHETIC\_SIM; not RF).}
\begin{tabular}{rcccc}\toprule
seed & mean demand (Mbps) & p95 demand (Mbps) & Jain fairness & radio evidence \\\midrule
1 & 25.4966 & 47.6709 & 0.7602 & synthetic\_fixture \\
2 & 24.9347 & 47.2698 & 0.7571 & synthetic\_fixture \\
7 & 25.0176 & 47.7018 & 0.7558 & synthetic\_fixture \\
42 & 26.0733 & 47.726 & 0.7749 & synthetic\_fixture \\
\bottomrule\end{tabular}
\end{table}
% sha 7d4aad3ce0e9a969


\section{Limitations}
\label{sec:limitations}
This is a digital synthetic-fixture study~\cite{itur_m2160}.
It is not RF sounding, not independent human reproduction, not a community deployment, not operator KPI, and not a University of Oulu affiliation claim.
Held-out physical campaigns remain \texttt{MEASUREMENT\_PENDING}.
The Gary bearer-stress mapping is a protocol assumption.
Seeded demand variation does not simulate channel fading; radio YAML is constant.
Comparative site inclusion scores are barrier-list composites, not field surveys.
\texttt{DIGITAL\_RESULTS\_VALIDATED} here means the digital path reproduced and the failure-case findings were extracted from JSON --- not that a conference paper is camera-ready or submitted.

\bibliographystyle{plain}
\bibliography{bibliography}
\end{document}
